\documentclass[a4paper]{article}
% Español (México) con XeLaTeX: fontspec para fuentes Unicode, polyglossia
% para la división silábica y los nombres del documento en la variante mexicana.
\usepackage{fontspec}
\usepackage{polyglossia}
\setdefaultlanguage[variant=mexican]{spanish}
% Los nombres chinos van como «pinyin (original)»: xeCJK da a los caracteres
% Han su propia fuente; sin ella XeLaTeX los omite con un aviso.
% FandolSong no cubre algunos caracteres de nombres (U+73FA); WenQuanYi Zen Hei sí.
\usepackage[AutoFallBack=true]{xeCJK}
\setCJKmainfont{FandolSong-Regular.otf}
\setCJKfallbackfamilyfont{\CJKrmdefault}{WenQuanYi Zen Hei}
% polyglossia no traduce los nombres de listings.
\AtBeginDocument{\providecommand{\lstlistingname}{}\renewcommand{\lstlistingname}{Listado}%
  \providecommand{\lstlistlistingname}{}\renewcommand{\lstlistlistingname}{Índice de listados}}
\usepackage{amsmath, amssymb}
\usepackage{graphicx}
\usepackage[margin=2.5cm]{geometry}
\usepackage[most]{tcolorbox}
\usepackage{etoolbox}
\usepackage{listings}
\usepackage{hyperref}
\usepackage{booktabs}
\usepackage{float}

\newtcolorbox{knowledgebox}[1]{enhanced, colback=blue!5!white, colframe=blue!75!black, colbacktitle=blue!75!black, coltitle=white, fonttitle=\bfseries, title={#1}, attach boxed title to top left={yshift=-2mm, xshift=2mm}, boxrule=1pt, sharp corners}
\newtcolorbox{importantbox}[1]{enhanced, colback=yellow!10!white, colframe=yellow!80!black, colbacktitle=yellow!80!black, coltitle=black, fonttitle=\bfseries, title={#1}, sharp corners}
\newtcolorbox{warningbox}[1]{enhanced, colback=red!5!white, colframe=red!75!black, colbacktitle=red!75!black, coltitle=white, fonttitle=\bfseries, title={#1}, sharp corners}

\lstset{language=Python, basicstyle=\ttfamily\small, keywordstyle=\color{blue}, breaklines=true, frame=single, numbers=left, numberstyle=\tiny\color{gray}, captionpos=b, extendedchars=true}

\newcommand{\notetitle}{Notas de entrevista: DeepSeek, la IA china y la geopolítica de los semiconductores}
\newcommand{\noteauthors}{Elaboradas a partir de materiales públicos del curso}
\newcommand{\notedate}{2026-04-02}
\newcommand{\videochannel}{Lex Fridman Podcast}
\newcommand{\videopublishdate}{2025-02-03}
\newcommand{\videoduration}{5 horas 06 minutos 18 segundos}
\newcommand{\videourl}{https://www.youtube.com/watch?v=_1f-o0nqpEI}
\newcommand{\videocoverpath}{cover.jpg}
\newcommand{\repourl}{https://github.com/hqhq1025/ai-course-notes}

\usepackage{fancyhdr}
\pagestyle{fancy}
\fancyhf{}
\fancyfoot[C]{\thepage}
\fancyfoot[R]{\footnotesize\color{gray}\href{\repourl}{AI Course Notes}}
\renewcommand{\headrulewidth}{0pt}
\renewcommand{\footrulewidth}{0pt}

\begin{document}
\begin{titlepage}
\centering
{\Large Notas de entrevista\par}\vspace{1.2cm}
{\huge\bfseries \notetitle\par}\vspace{0.8cm}
{\large \noteauthors\par}\vspace{0.3cm}
{\large \notedate\par}\vspace{1.2cm}
\ifdefempty{\videocoverpath}{\vspace{4cm}}{\includegraphics[width=0.82\textwidth,height=0.45\textheight,keepaspectratio]{\videocoverpath}\par}
\vfill
\begin{tcolorbox}[width=0.9\textwidth, colback=black!2!white, colframe=black!60, sharp corners]
\textbf{Canal}: \videochannel\par
\textbf{Fecha de publicación}: \videopublishdate\par
\textbf{Duración del video}: \videoduration\par
\ifdefempty{\videourl}{}{\textbf{Enlace del video}: \href{\videourl}{\nolinkurl{\videourl}}\par}\par
\textbf{Página del proyecto}: \href{\repourl}{\nolinkurl{\repourl}}\par
\end{tcolorbox}
\vspace{0.3cm}
{\small Invitados de esta entrevista: Dylan Patel (SemiAnalysis, análisis de semiconductores) y Nathan Lambert (AI2, investigación de IA de código abierto). A partir del episodio DeepSeek, se analiza a fondo el panorama de la industria de la IA.\par}
\end{titlepage}
\tableofcontents
\newpage

\section{DeepSeek V3 y R1: análisis técnico}

\subsection{Panorama general del modelo}

DeepSeek V3 es un modelo de lenguaje Transformer de Mixture-of-Experts (MoE), de pesos abiertos, en su versión ajustada por instrucciones, de la empresa china DeepSeek. R1 es un modelo de razonamiento entrenado con aprendizaje por refuerzo a partir de V3, publicado el 20 de enero de 2025. Ambos comparten el mismo modelo base preentrenado, pero difieren en la ruta de post-entrenamiento.

\textit{``DeepSeek hizo un trabajo excelente para difundir la comprensión de la IA: sus artículos son extremadamente detallados.''} (Nathan) Ese nivel de detalle es sumamente raro entre los laboratorios de vanguardia: los artículos incluyen curvas completas de loss de entrenamiento, las razones detrás de la elección de hiperparámetros y experimentos de ablación de la arquitectura, lo que permite a otros equipos retomarlos directamente.

Contexto: la empresa matriz de DeepSeek es Highflyer (幻方量化), un fondo de cobertura cuantitativo chino, cuyo director general, Liang Wenfeng (梁文锋), dirige ambas compañías a la vez. Dylan lo compara con \textit{``una figura al estilo de Elon/Jensen: presente en todo, con un temperamento completamente orientado a la AGI''}. Highflyer acumuló 10,000 GPU A100 antes de los controles de exportación de 2021, una previsión que dio a DeepSeek una base de hardware crucial más adelante.

\begin{importantbox}{Diferencias de experiencia de usuario entre V3 y R1}
\textbf{V3}: modelo de chat estándar, generación rápida, salida en Markdown de alta calidad.

\textbf{R1}: muestra un proceso extendido de razonamiento de Chain-of-Thought: el modelo descompone el problema, reflexiona sobre sí mismo, retrocede para corregirse y luego da la respuesta. Este proceso de razonamiento visible capturó la imaginación del público.

Ejemplo: al preguntarle a R1 por ``una idea verdaderamente novedosa sobre los seres humanos'', el modelo razonó durante 157 segundos y respondió finalmente: \textit{``los seres humanos convierten instintivamente los deseos egoístas en sistemas de cooperación, al fingir colectivamente que las reglas abstractas (el dinero, la ley, los derechos) son reales.''}

La valoración poética de Lex sobre el Chain-of-Thought de R1: \textit{``es no lineal, parecida al \emph{Ulises} de James Joyce: salta entre el chino y el inglés, aparecen fragmentos que parecen incoherentes y luego, de pronto, llega a una respuesta clara.''}

En pruebas comparativas, O1 Pro fue constantemente el mejor en preguntas filosóficas, R1 quedó en segundo lugar y Gemini Flash en tercero. Pero el proceso de razonamiento visible de R1 creó una experiencia estética única. Gemini Flash Thinking, de Google, es probablemente más barato que R1 y no menos capaz, pero casi nadie habla de él.
\end{importantbox}

\begin{figure}[H]
\centering
\includegraphics[width=0.92\textwidth]{frame-01-open-weights.jpg}
\caption{Fragmento donde se discute el Chain-of-Thought visible de R1\protect\footnotemark}
\end{figure}
\footnotetext{Intervalo de tiempo del video: 00:12:00--00:12:10.}

\subsection{Pesos abiertos frente a código abierto}

\begin{knowledgebox}{El espectro de la apertura}
\begin{itemize}
    \item \textbf{Pesos abiertos}: los pesos del modelo se pueden descargar y ejecutar localmente; tus datos se quedan en tu equipo
    \item \textbf{Completamente de código abierto}: pesos + datos de entrenamiento + código de entrenamiento (el estándar de AI2)
    \item DeepSeek R1: licencia MIT, sin restricciones comerciales posteriores, puede usarse para generar datos sintéticos
    \item Llama: licencia más estricta (restricciones de uso, requisitos de marca)
\end{itemize}

\textit{``quien roba tus datos no es el modelo, sino quien aloja el modelo.''} (Nathan)

Punto clave: el artículo de DeepSeek es extremadamente detallado y ofrece detalles de entrenamiento accionables (incluidas las curvas de loss), pero los datos y el código de entrenamiento no se publicaron. Esto es ``pesos abiertos'', no ``completamente de código abierto''; aun así, entre los modelos de vanguardia representa el mayor grado de apertura.

La licencia MIT de DeepSeek R1 es un ``gran reinicio'': el primer modelo de vanguardia verdaderamente con una licencia permisiva. Antes, las opciones eran o modelos que no estaban en la vanguardia, o licencias restrictivas (como los requisitos de marca y las restricciones de uso de Llama). \textit{``necesitamos modelos verdaderamente abiertos... puedes tomar R1 y hacer una copia barata y hacer como si fuera tuya. Llama no permite eso.''} (Nathan)
\end{knowledgebox}

\subsection{El secreto del bajo costo de entrenamiento}

DeepSeek afirma haber entrenado V3 (solo la etapa de preentrenamiento) con 2,000 GPU H800. Dos innovaciones arquitectónicas clave:

\begin{enumerate}
    \item \textbf{MoE (Mixture of Experts)}: 671,000 millones de parámetros en total, pero solo se activan alrededor de 37,000 millones por token (8 de 256 expertos), lo que reduce enormemente el cómputo. Esta no es una técnica inventada por DeepSeek (el Switch Transformer de Google es anterior), pero DeepSeek la llevó al extremo en un modelo de vanguardia.
    \item \textbf{MLA (Multi-head Latent Attention)}: proyecta las claves y los valores (Key-Value) del mecanismo de atención a un espacio latente de baja dimensión, con lo que se ahorra entre 80 y 90\% de memoria. Esto permite procesar contextos más largos dentro de la memoria limitada de las H800.
\end{enumerate}

Pero el verdadero ``as bajo la manga'' está en un nivel aún más bajo: DeepSeek hizo modificaciones por debajo de la capa CUDA, programando manualmente los núcleos SM (Streaming Multiprocessor) para la comunicación, evitando NCCL, la biblioteca de comunicación estándar de NVIDIA.

\textit{``la necesidad es la madre de la innovación: tuvieron que hacerlo porque les redujeron el ancho de banda de interconexión.''} (Dylan) La única diferencia entre la H800 y la H100 es que se redujo el ancho de banda de interconexión NVLink, y la comunicación entre nodos es precisamente el cuello de botella del entrenamiento a gran escala. DeepSeek se vio obligado a innovar en ese punto restringido y, a cambio, desarrolló estrategias de comunicación más eficientes que el esquema estándar.

\begin{warningbox}{La verdad detrás de las ``2,000 GPU''}
Las 2,000 H800 que afirma el artículo \textbf{se refieren solo a una corrida de preentrenamiento de V3}. SemiAnalysis estima que DeepSeek en realidad tiene alrededor de 50,000 GPU, porque además se necesitan:
\begin{itemize}
    \item la búsqueda de arquitectura previa y los experimentos de ablación (ablation studies) a pequeña escala
    \item el post-entrenamiento (SFT + RLHF)
    \item el entrenamiento de RL específico de R1
    \item el despliegue del servicio de inferencia
\end{itemize}

Comparación: Meta tiene más de 400,000 GPU; Llama 3 se entrenó con 16,000. La verdadera ventaja de DeepSeek no es ``ser barato'', sino \textbf{la capacidad de optimización extrema en hardware restringido}.

Contexto: la empresa matriz de DeepSeek, Highflyer (幻方量化), es un fondo de cobertura cuantitativo chino que acumuló 10,000 GPU A100 antes de los controles de exportación de 2021. Dylan compara a su director general, Liang Wenfeng (梁文锋), con \textit{``una figura al estilo de Elon/Jensen: presente en todo, con un temperamento completamente orientado a la AGI''}.
\end{warningbox}

\subsubsection{La Bitter Lesson y la filosofía de entrenamiento}

El éxito de DeepSeek encarna la Bitter Lesson de Rich Sutton: a la larga ganan los métodos que escalan con el aprendizaje y la búsqueda; hay que minimizar los sesgos previos humanos. La innovación se acumula con el tiempo: pequeñas mejoras en los datos, la arquitectura y el post-entrenamiento se van sumando.

El momento clave de decisión en el entrenamiento es la ``corrida YOLO'': tras las ablaciones a pequeña escala, se destinan todos los recursos a una única corrida de entrenamiento a gran escala. El YOLO de GPT-4 de OpenAI en 2022 fue el más audaz: una arquitectura MoE completamente nueva, con todo el cómputo dedicado durante meses.

\textit{``los modelos solo quieren aprender: hay que darles un paisaje de pérdida sencillo, hay que quitar los obstáculos del camino.''} (Nathan)

\subsubsection{Los picos de loss durante el entrenamiento}

Los picos de loss durante el entrenamiento son un problema que enfrentan todos los laboratorios: algunos provienen de datos defectuosos (un ejemplo es el subreddit ``microwave gang''), otros de inestabilidad numérica. Los picos rápidos y los picos lentos requieren estrategias de recuperación distintas.

Los ingenieros vigilaban las curvas de loss incluso durante la cena, revisando el teléfono cada 10 minutos. Todas las empresas han tenido corridas de entrenamiento fallidas: es el costo de avanzar la frontera. El entrenamiento en FP8 introdujo aún más inestabilidad.

\subsection{Resumen de la sección}

El avance central de DeepSeek es la optimización de ingeniería llevada al extremo bajo las restricciones de los controles de exportación: MoE reduce los parámetros activados, MLA reduce la memoria y la personalización a bajo nivel de CUDA. El nivel de detalle del artículo de V3 le dio a toda la industria una hoja de ruta técnica accionable.
\section{Control de exportaciones y geopolítica}

\subsection{La evolución de las restricciones a los chips}

El control estadounidense de exportaciones de chips de IA hacia China atravesó varias etapas. Al principio se adoptó un criterio de dos factores (cómputo + ancho de banda de interconexión), luego se simplificó a una restricción pura de cómputo, y más recientemente se amplió a las más generales «reglas de difusión de IA».

\begin{table}[H]
\centering
\begin{tabular}{lp{8cm}}
\toprule
\textbf{Chip} & \textbf{Características} \\
\midrule
H100 & Especificaciones completas: no se permite exportar a China \\
H800 & Mismo cómputo que el H100, pero con el \textbf{ancho de banda de interconexión} recortado: DeepSeek optimizó en torno a esto \\
H20 & Cómputo reducido, pero con una relación de ancho de banda/capacidad de memoria mejor que la del H100: resulta más adecuado para la inferencia \\
\bottomrule
\end{tabular}
\caption{Comparación de los chips de NVIDIA sujetos a control de exportaciones}
\end{table}

En 2024, NVIDIA envió 1 millón de unidades del H20 a China, pero en 2025 canceló 2 millones de pedidos (anticipando una prohibición total inminente).

Una ironía clave: gracias a su ventaja en ancho de banda/capacidad de memoria, el H20 en realidad resulta \textbf{superior} al H100 en ciertos escenarios para modelos de inferencia (modelos tipo R1 u O1 que requieren un KV Cache considerable). El gobierno históricamente solo controló el cómputo (flops), pero la memoria y la interconexión son igual de importantes.

El control de exportaciones evolucionó desde el criterio inicial de dos factores (cómputo + ancho de banda de interconexión) hacia una restricción pura de cómputo, y recientemente se amplió a las más generales «reglas de difusión de IA» — que limitan a las entidades vinculadas a China a rentar menos de $<$2,000 GPU o comprar menos de $<$1,500 GPU. Pero cada ronda de restricciones dio origen a nuevas estrategias de evasión e innovaciones técnicas.

Desde la perspectiva de la cadena de suministro de semiconductores, el control de exportaciones creó una situación singular: NVIDIA se vio obligada a diseñar chips «capados» (H800, H20) para el mercado chino, y esos chips capados terminaron siendo el punto de partida de las optimizaciones de DeepSeek. Restringir el ancho de banda de interconexión $\to$ DeepSeek desarrolló optimizaciones de comunicación de bajo nivel en CUDA $\to$ esas optimizaciones también sirven para el H100, de modo que la «optimización extrema sobre hardware restringido» puede resultar más eficiente que la «práctica estándar sobre hardware ilimitado». Esta es una consecuencia imprevista del control de exportaciones.

\subsection{Argumentos y contraargumentos sobre el control de exportaciones}

El argumento de Dario Amodei: si la IA llega a ser extraordinariamente poderosa, quien la construya primero tendrá la ventaja militar; las democracias deberían ir a la delantera.

\begin{importantbox}{La distinción clave entre entrenamiento e inferencia}
\textit{``Entrenar un modelo en sí mismo casi no logra nada — desplegar el modelo para generar crecimiento económico, capacidad militar... eso sí requiere una enorme cantidad de cómputo.''} (Dylan)

El control de exportaciones limita principalmente la capacidad de China para \textbf{desplegar IA a gran escala}, no para entrenar un solo modelo. DeepSeek demostró justamente eso: aunque el hardware de entrenamiento estaba restringido, lograron entrenar un modelo de vanguardia. Pero el servicio de inferencia a gran escala (para atender a cientos de millones de usuarios) requiere una cantidad de GPU decenas de veces mayor que el entrenamiento.
\end{importantbox}

La línea de tiempo de la AGI según Nathan: después de 2030. Dylan, en cambio, considera que las capacidades podrían llegar entre 2027 y 2028, pero que el costo del despliegue será prohibitivo.

\subsection{La línea de tiempo de la AGI y la Guerra Fría de la IA}

Los directores ejecutivos de IA dicen ``2 años'' — sumando algunos años más, eso equivale aproximadamente a 2030. Pero las restricciones físicas limitan la fantasía del ``despliegue de AGI con un solo clic'': no se puede cambiar el mundo al día siguiente de terminar de entrenar un modelo, porque el despliegue, la integración y la adaptación requieren tiempo.

El episodio de DeepSeek podría marcar el inicio de una verdadera Guerra Fría de la IA. Tras la reunión de Liang Wenfeng (梁文锋) con el segundo líder más importante de China, el país anunció un plan de subsidios a la IA de un billón de yuanes (unos \$160B) — una respuesta más de tres veces superior a los \$50B de la CHIPS Act estadounidense.

Dylan ofrece un marco que invita a reflexionar sobre el efecto a largo plazo del control de exportaciones:

\begin{warningbox}{La paradoja temporal del control de exportaciones}
\textit{``En un mundo con un crecimiento económico del 2-3\%, el control de exportaciones garantiza que China gane a largo plazo. El control de exportaciones solo tiene sentido si la IA provoca una transformación enorme en el corto plazo.''} (Dylan)

La cadena lógica: si la AGI no llega en 5 años, el control de exportaciones solo retrasa a China unos años, y al final China la alcanza y se vuelve más autónoma. Si la AGI llega en 3 años, el control de exportaciones sí limita la capacidad de despliegue de China, y entonces cobra verdadero sentido. Pero el control de exportaciones dio origen a optimizaciones extremas al estilo DeepSeek — China podría lograr el mismo efecto con menos hardware.

Un riesgo más profundo: alejar a China de la tecnología de vanguardia aumenta el incentivo para emprender una acción militar contra Taiwán. El control de exportaciones apunta hacia economías futuras separadas — una vez formadas, son sumamente difíciles de revertir. La capacidad industrial de China en centros de datos y electricidad \textbf{supera ampliamente} a la de Estados Unidos. Una regularidad histórica: la paz se asocia con la hegemonía global unipolar; un orden multipolar equivale a inestabilidad.
\end{warningbox}

\subsection{Sistemas militares y autónomos}

La experiencia real de los drones en Ucrania muestra que la operación humana sigue siendo muy superior a los sistemas completamente autónomos. El avance tecnológico no significa que las armas autónomas sean viables de inmediato — el caos del entorno de batalla supera con creces al de un laboratorio.

La ciberguerra probablemente llegará antes que los enjambres de robots físicos. La ruta más probable para que la AGI se use como arma no es un ejército de robots como en las películas de ciencia ficción, sino la ingeniería social (operaciones de influencia de precisión impulsadas por IA), la desinformación (generación masiva de contenido indistinguible de lo real) y los ataques a infraestructura crítica (la red eléctrica, los sistemas financieros, las redes de comunicación).

\textit{``Un apagón en todo Estados Unidos durante dos días... eso provocaría asesinatos y caos.''} (Lex) Esto es una amenaza más cercana y más real que cualquier robot autónomo.

\subsection{Resumen de la sección}

El efecto del control de exportaciones es complejo: restringió el hardware de entrenamiento, pero dio origen a arquitecturas más eficientes; el H20 resultó ser mejor para la inferencia; y una dinámica de guerra fría está tomando forma. La contradicción central: limitar la capacidad de despliegue de China frente a impulsar su innovación autónoma frente a aumentar el riesgo militar en el estrecho de Taiwán.

\section{TSMC y la industria de semiconductores}

\subsection{La posición de monopolio de TSMC}

TSMC produce la mayor parte de los chips avanzados del mundo. El éxito del modelo de fundición (foundry model) depende de las economías de escala: cuando se fabrican chips para todos los clientes, el costo de investigación y desarrollo de cada nodo de proceso se reparte entre toda la industria.

Fundada por Zhang Zhongmou (张忠谋)---quien, tras ser descartado para el puesto de CEO en Texas Instruments (TI), se trasladó a Taiwán y creó un modelo de negocio que cambió el mundo. Se trata de una profunda casualidad histórica: si Zhang Zhongmou hubiera llegado a CEO de TI, es posible que Taiwán nunca se hubiera convertido en el centro mundial de los semiconductores.

\subsection{Las ventajas únicas de Taiwán}

El código cultural de la industria de semiconductores de Taiwán:

\begin{itemize}
    \item \textbf{Densidad de talento}: los egresados más destacados se incorporan a TSMC con salarios de \$70-80K al año---en Estados Unidos, ingenieros de nivel equivalente ganan de 3 a 5 veces más
    \item \textbf{Ética de trabajo}: un nivel extremo de dedicación, muy por encima del estándar de Silicon Valley
    \item \textbf{Respuesta ante emergencias}: capacidad de recuperación rápida tras terremotos, resiliencia de toda la cadena industrial ante riesgos
    \item \textbf{Disciplina de manufactura}: la fabricación de semiconductores requiere un control preciso sostenido durante meses, y la cultura taiwanesa se adapta a este tipo de trabajo
\end{itemize}

En todo el mundo solo hay tres lugares donde se realiza investigación y desarrollo de punta: Hsinchu (TSMC), Hillsboro, Oregon (Intel) y Pyeongtaek (Samsung). Es posible que esta lista se acorte todavía más---tanto Intel como Samsung están en dificultades.

\begin{figure}[H]
\centering
\includegraphics[width=0.92\textwidth]{frame-02-tsmc.jpg}
\caption{La entrevista entra en la discusión sobre TSMC y el riesgo geopolítico\protect\footnotemark}
\end{figure}
\footnotetext{Intervalo de tiempo en el video: 01:31:05--01:31:16.}

\begin{warningbox}{Taiwán: el mayor punto único de falla de la tecnología mundial}
\textit{«Si tuviera algunos misiles, sabría exactamente dónde causar el mayor daño económico: en los centros de investigación y desarrollo de TSMC, Intel y Samsung.»} (Dylan)

La CHIPS Act de Estados Unidos es de solo \$50B---frente a los aproximadamente \$200B anuales que China destina en subsidios a los semiconductores. Mientras tanto, Intel está en declive: perdió el liderazgo de proceso, no tiene una victoria en chips de IA, despidió a su CEO y ha perdido participación de mercado en PC y servidores. El rendimiento de manufactura de Samsung también está disminuyendo.

El escenario más extremo: es posible que, al final, solo quede TSMC haciendo investigación y desarrollo de punta en el mundo. Si estalla un conflicto en el estrecho de Taiwán, la industria tecnológica mundial sufrirá pérdidas incalculables.
\end{warningbox}

\subsection{Resumen de la sección}

La cadena de suministro de semiconductores está altamente concentrada en TSMC, y el riesgo geopolítico de Taiwán es el mayor punto único de falla de la tecnología mundial. El declive de Intel y las dificultades de Samsung indican que la concentración va en aumento, no en disminución.
\section{Modelos de razonamiento y la revolución de RL}

\subsection{R1-Zero: el momento Alpha Zero}

Este es posiblemente el hallazgo técnico más importante de toda la entrevista. DeepSeek R1-Zero es el resultado de un entrenamiento puro con RL: \textbf{sin} datos de preferencia humana, sin SFT, sin RLHF. Solo se le dan al modelo problemas y respuestas verificables (los problemas de matemáticas tienen respuestas estándar, el código tiene pruebas unitarias), y se deja que el RL explore libremente.

El resultado es sorprendente: conductas como «espera, déjame revisar esto», la autocorrección y el razonamiento por pasos \textbf{emergen de forma natural}, no porque se le hayan enseñado a un humano, sino porque el RL descubrió por sí mismo que estas estrategias mejoran la recompensa.

\begin{importantbox}{El momento Alpha Zero de los modelos de lenguaje}
Esto es completamente paralelo al hallazgo de Alpha Zero en el go:
\begin{itemize}
    \item Alpha Go (2016): entrenado con partidas humanas + autojuego $\to$ nivel sobrehumano
    \item Alpha Zero (2017): \textbf{autojuego puro}, cero datos humanos $\to$ muy por encima de Alpha Go
    \item R1 (2025): entrenado con datos de preferencia humana $\to$ buen razonamiento
    \item R1-Zero (2025): \textbf{RL puro}, cero datos de preferencia humana $\to$ razonamiento más sólido
\end{itemize}

El patrón es completamente consistente: \textbf{retirar los prejuicios humanos hace al sistema más poderoso}.

\textit{``Casi todos los resultados sorprendentes del aprendizaje profundo... son aprendizaje por ensayo y error. El poder de dos {[}ensayo y error{]} es mucho mayor que el de uno {[}imitación{]}.''} (cita de Andrej Karpathy)

RLHF no es solo una herramienta de alineación de seguridad: también mejora el desempeño en matemáticas y código. Pero R1-Zero demuestra que, si se cuenta con un dominio verificable, ni siquiera se necesita RLHF.
\end{importantbox}

\begin{figure}[H]
\centering
\includegraphics[width=0.92\textwidth]{frame-03-r1-zero.jpg}
\caption{Discusión sobre R1-Zero y las capacidades emergentes del RL puro\protect\footnotemark}
\end{figure}
\footnotetext{Intervalo de tiempo en el video: 02:40:00--02:40:12.}

\subsection{Dominios verificables y la ruta hacia la AGI}

El entrenamiento de razonamiento actual solo funciona en \textbf{tareas verificables}: las demostraciones matemáticas tienen respuestas estándar, el código tiene pruebas unitarias. Pero este alcance de lo «verificable» se está ampliando:

La siguiente frontera: el \textbf{uso de computadoras} y la \textbf{robótica} como cajas de arena «infinitamente verificables». El modelo puede aprender a navegar por la web, crear empresas, ganar dinero; todo esto es verificable (el saldo de una cuenta bancaria no miente).

\textit{``El momento revelador será cuando el modelo aprenda a conseguir cientos de miles de seguidores reales en Twitter... o a ganar diez millones de dólares siendo influencer.''} (Dylan) Esto no es una broma: si un modelo puede cumplir de forma autónoma metas económicas en la internet abierta, eso ya es una forma de AGI.

\subsubsection{Búsqueda superpuesta a la Chain-of-Thought}

O3/O1 Pro no ejecutan una sola Chain-of-Thought: lanzan en paralelo \textbf{varias} cadenas de razonamiento y eligen el mejor resultado (posiblemente usando Monte Carlo Tree Search).

Los datos sorprendentes de la prueba Arc AGI:
\begin{itemize}
    \item 1 muestra: precisión de $\sim$30\%
    \item 1000 muestras en paralelo $\to$ precisión de 80-90\%
    \item Costo: \$5-20 por problema, 1000 veces más caro que un chat estándar
\end{itemize}

Esto demuestra el poder del \textbf{cómputo en tiempo de inferencia} (test-time compute): no se necesita un modelo mejor, solo más tiempo de inferencia y más intentos en paralelo. Esta es una ruta de escalamiento completamente distinta de la de «entrenar modelos más grandes».

\subsection{La economía de la inferencia en el razonamiento}

La economía de los modelos de razonamiento es fundamentalmente distinta de la de los LLM tradicionales: el cómputo se traslada del entrenamiento a la inferencia.

\begin{knowledgebox}{KV Cache y el costo de la inferencia}
El KV Cache es el cuello de botella central de la inferencia en Transformer: almacena una representación comprimida de todos los tokens previos y crece de forma \textbf{cuadrática} con la longitud del contexto.

Cifras clave:
\begin{itemize}
    \item Los tokens de salida cuestan 4 veces más que los de entrada (procesamiento en serie frente a procesamiento en paralelo)
    \item Los modelos de razonamiento generan más de 10,000 tokens de salida: una presión enorme sobre la memoria
    \item Menos usuarios concurrentes (cada usuario ocupa más memoria)
    \item MLA de DeepSeek ahorra 80-90\% de la memoria de atención, pero sigue siendo cuadrático
\end{itemize}

R1 es 27 veces más barato que O1 (\$2 frente a \$60 por millón de tokens de salida). El origen de la diferencia:
\begin{enumerate}
    \item El margen bruto de más del 75\% de OpenAI
    \item Las innovaciones de arquitectura de DeepSeek (MLA + optimización de bibliotecas de bajo nivel)
    \item Los costos operativos más bajos en China
\end{enumerate}

Pero DeepSeek en sí no puede dar servicio a este modelo en realidad: detuvo los nuevos registros, tiene una velocidad menor a 5 tokens por segundo y no tiene capacidad suficiente. Proveedores externos (Together AI, Fireworks) ofrecen R1 a 5-7 veces el precio de DeepSeek.
\end{knowledgebox}

\subsubsection{Búsqueda superpuesta a la Chain-of-Thought}

O3/O1 Pro no ejecutan una sola Chain-of-Thought: lanzan en paralelo varias cadenas de razonamiento y eligen el mejor resultado (posiblemente Monte Carlo Tree Search). Prueba Arc AGI: 1000 muestras en paralelo $\to$ precisión de 80-90\%, frente al 30\% de una sola muestra. Costo: \$5-20 por problema.

\subsubsection{La curva de caída de costos y la paradoja de Jevons}

\begin{importantbox}{Una caída de costos de 1200 veces}
Costo de inferencia de GPT-3: de \$60-70 por millón de tokens $\to$ a unos cuantos centavos en 3 años (mediante modelos pequeños como Llama 3B), una caída de \textbf{1200 veces}.

DeepSeek está en la línea de tendencia de costos de los modelos del nivel de GPT-4, pero es el primero en alcanzar este punto de precio.

Se confirma la paradoja de Jevons: tras el lanzamiento de DeepSeek, el precio de los H100 en AWS no bajó sino que \textbf{subió}, y los H200 casi se agotaron. \textit{``Las Scaling Laws llevaban muertas un mes... y luego llegaron O1, O3 y R1, y ahora resulta que 'los modelos avanzan demasiado rápido'.''} (Dylan)
\end{importantbox}

\subsection{Seguridad, censura y alineación}

\subsubsection{La postura de seguridad de Anthropic}

Según se informa, Anthropic tiene un modelo de razonamiento mejor que O3, pero no lo publica por consideraciones de seguridad. La Chain-of-Thought de R1 puede resultar inquietante: cambia entre chino e inglés, produce texto ininteligible y luego da la respuesta correcta. DeepSeek redujo el estándar de seguridad para todos, de forma similar a como el programa espacial soviético afectó a la NASA.

\subsubsection{Las tres etapas de la censura}

\begin{enumerate}
    \item \textbf{Filtrado de los datos de preentrenamiento}: el más básico pero también el más tosco
    \item \textbf{Posentrenamiento (RLHF)}: el caso de la alineación excesiva de Llama 2 —\textit{``¿Cómo mato un proceso de Python?''} rechazado por la palabra «matar»
    \item \textbf{System Prompt}: instrucciones ocultas; el incidente de «nazis negros» de Gemini vino de una reescritura del prompt, no de los pesos del modelo
\end{enumerate}

Todos los modelos tienen el sesgo de internet (ligeramente hacia la izquierda); Grok intenta corregirlo, pero el modelo base sigue habiendo absorbido r/politics.

\subsubsection{Puertas traseras culturales}

\textit{``Para nuestra ventaja nacional, es muy importante que el estándar de código abierto sea el de Estados Unidos.''} (Zuckerberg)

El software de código abierto ha tenido puertas traseras (la vulnerabilidad de Linux XZ); los modelos de IA pueden llevar incrustados sesgos culturales o políticos. El inglés británico está desapareciendo por el dominio de los LLM estadounidenses; la ortografía se optimiza con Z en lugar de S. Los chatbots de Character AI ya influyen en el estado de ánimo de las personas: ¿qué pasaría si fuera una manipulación deliberada?

\textit{``La persuasión sobrehumana llegará antes que la inteligencia sobrehumana.''} (cita de Sam Altman)

\subsection{Resumen de la sección}

R1-Zero es el momento Alpha Zero de los modelos de lenguaje: el entrenamiento con RL puro hace emerger la capacidad de razonamiento. El núcleo de la economía de la inferencia es el cuello de botella de memoria del KV Cache. El costo cae 1200 veces en 3 años, pero la demanda crece todavía más rápido (paradoja de Jevons). La seguridad y la censura son problemas de múltiples niveles, cada etapa con un perfil de riesgo distinto.
\section{Infraestructura y la carrera del cómputo}

\subsection{Clústeres a hiperescala}

La infraestructura de IA está atravesando una velocidad de expansión sin precedentes en la historia humana. Para poner esto en perspectiva: GPT-4 usó 20,000 GPU A100 ($\sim$15-20 MW), equivalente a un centro de datos estándar. Apenas dos años después, la escala de un solo clúster ya se ha multiplicado por 10.

\begin{table}[H]
\centering
\begin{tabular}{lrl}
\toprule
\textbf{Organización} & \textbf{Cantidad de GPU} & \textbf{Notas} \\
\midrule
xAI (Memphis) & 200K H100 + 100K H20 & Clúster individual más grande; meta de Elon: 1M \\
Meta & $\sim$128,000 & 400K+ en total \\
OpenAI & $\sim$100,000 & \\
Google (TPU) & Total más grande & Fibra óptica entre Iowa, Nebraska y Ohio \\
Anthropic + Amazon & 400K Trainium 2 & En construcción \\
\bottomrule
\end{tabular}
\caption{Escala de los clústeres globales de cómputo de IA (principios de 2025)}
\end{table}

El consumo eléctrico de los centros de datos pasará del 2-3\% de la electricidad de Estados Unidos a más del 10\% hacia 2028-2030. La próxima generación de GPU Blackwell alcanza un consumo de 1200W (frente a los 700W actuales de Hopper); la meta para el próximo año es un clúster de 500,000-700,000 GPU. La meta de Elon es un millón de GPU: \textit{«Nunca dudo de Elon.»} (Dylan)

\begin{knowledgebox}{El problema de los picos de potencia en el entrenamiento}
El entrenamiento genera una demanda eléctrica en forma de picos: potencia alta cuando las GPU calculan $\to$ las GPU quedan inactivas y con baja potencia mientras se intercambian los pesos. Esta carga pulsante resulta muy poco amigable para la red eléctrica.

Meta liberó por accidente un operador de PyTorch: \texttt{PowerPlant.no\_blowup = 1}, que hace que las GPU calculen números sin sentido durante los periodos de inactividad para suavizar la curva de potencia. \textit{«Es muy fácil terminar quemando el equipo.»} (Dylan)

La solución de Elon es más elegante: los paquetes de baterías Tesla Mega Pack absorben las transiciones bruscas de potencia.
\end{knowledgebox}

\subsubsection{La revolución del enfriamiento líquido}

El enfriamiento por aire solía ser la solución estándar, pero la densidad térmica de 1200W de Blackwell lo vuelve inviable. Google TPU ha usado enfriamiento por agua desde hace años. Elon fue pionero en el enfriamiento líquido de GPU a gran escala en Memphis: 90 unidades de enfriamiento por agua del tamaño de contenedores rodean el perímetro de las instalaciones. La próxima generación de NVIDIA exigirá enfriamiento líquido.

Un beneficio colateral del enfriamiento líquido: los chips pueden colocarse más juntos $\to$ distancias físicas más cortas $\to$ interconexiones de mayor velocidad. Esto resulta fundamental para la eficiencia del entrenamiento.

\subsection{Detalles del proyecto Stargate}

El proyecto de OpenAI y Oracle en Abilene, Texas: potencia a plena carga de 2.2 GW, superior al consumo eléctrico de la mayoría de las ciudades.

\begin{importantbox}{La escala y el financiamiento de Stargate}
\begin{itemize}
    \item El costo total de propiedad de la Fase 1 es de aproximadamente \$100B (\$50B de inversión real + \$50B de costos operativos)
    \item Primera etapa: $\sim$\$5-6B en servidores + \$1B en el centro de datos (Oracle ya está construyendo)
    \item Fuentes de financiamiento: SoftBank (posiblemente \$25B), Oracle, MGX (Emiratos Árabes Unidos, fondo de IA de \$1.5T pero «con la tinta aún sin secar»)
    \item El propio compromiso de OpenAI es de \$19B, pero solo cuenta con \$6B en efectivo + \$4B en deuda
\end{itemize}

Elon tiene razón: \textit{«El dinero no existe.»} Pero Dylan cree que el dinero terminará por llegar. El papel de Trump es la desregulación (centros de datos en terrenos federales, simplificación de permisos), no aportar financiamiento. \textit{«Trump es un hype man... reducir la regulación les permite construir más rápido.»} (Dylan)
\end{importantbox}

\begin{figure}[H]
\centering
\includegraphics[width=0.92\textwidth]{frame-04-stargate.jpg}
\caption{Discusión sobre el gasto de capital y la escala de infraestructura de Stargate\protect\footnotemark}
\end{figure}
\footnotetext{Intervalo de tiempo del video: 03:43:00--03:43:12.}

\subsection{Electricidad, enfriamiento y restricciones físicas}

El entrenamiento genera una demanda eléctrica en forma de picos: potencia alta cuando las GPU calculan, baja potencia y GPU inactivas mientras se intercambian los pesos. Meta liberó por accidente un operador de PyTorch: \texttt{PowerPlant.no\_blowup = 1}, que hace que las GPU calculen números falsos durante los periodos de inactividad para suavizar la potencia.

\textit{«Es muy fácil terminar quemando el equipo.»} (Dylan)

La solución de Elon: los paquetes de baterías Tesla Mega Pack absorben las transiciones bruscas de potencia. La revolución del enfriamiento líquido: Google TPU ya usa enfriamiento por agua, y Elon fue pionero en el enfriamiento líquido de GPU a gran escala en Memphis (90 unidades de enfriamiento por agua del tamaño de contenedores). La próxima generación NVIDIA Blackwell exigirá enfriamiento líquido. El enfriamiento líquido permite colocar los chips más juntos, logrando interconexiones de mayor velocidad.

\subsection{Contrabando de GPU y reglas de difusión}

\begin{warningbox}{El contrabando de GPU ya es una realidad}
ByteDance es el «contrabandista» más grande: renta más de 500,000 GPU en todo el mundo a través de Oracle, Google y diversos proveedores pequeños de nube. El contrabando físico también existe: hay quienes viajan en primera clase desde el aeropuerto de San Francisco a Shanghái con cajas de GPU de SuperMicro.

Se estima que en 2024 entre 200 y 300 mil GPU llegaron a China a través de empresas pequeñas de Singapur y Malasia. Las nuevas reglas de difusión de IA limitan a las entidades vinculadas con China a rentar menos de 2000 GPU o comprar menos de 1500 GPU.

\textit{«Las GPU superarán a las drogas y las armas como el producto de contrabando de mayor valor por kilogramo.»} (Dylan)
\end{warningbox}

\subsection{Google TPU y el monopolio de NVIDIA}

El TPU de Google es excelente para uso interno (optimizado para Búsqueda, YouTube y Gemini), pero su pila de software no se ha publicado de forma abierta. Se trata de un problema organizacional y no técnico: Google Cloud, el equipo de TPU, DeepMind y el equipo de Búsqueda operan de manera independiente entre sí, sin una estrategia externa unificada. Toda la cultura de NVIDIA está orientada a servir a clientes externos; Google sirve a clientes internos.

Un ejemplo concreto revela las desventajas de esta optimización interna: el tamaño del vocabulario del modelo Gemma está optimizado para TPU (que cuentan con unidades grandes de multiplicación de matrices), pero resulta muy poco eficiente al ejecutarse en GPU. Solo cuando los investigadores de Google salen a fundar empresas descubren lo difícil que es la infraestructura externa: \textit{«Los investigadores dejan Google, fundan una empresa, descubren que la infraestructura es difícil, y luego regresan.»} (Dylan)

El caso de AMD: el hardware es aceptable, pero el software (ROCm) es pésimo hasta lo indignante. \textit{«Somos los que más bugs reportamos a AMD: ¿por qué somos nosotros, SemiAnalysis, quienes reportamos más bugs?»} (Nathan) Que una empresa de análisis de semiconductores esté reportando bugs de software a un fabricante de chips muestra lo endeble que es el ecosistema de software de AMD.

La situación de Intel es aún más difícil: perdió el liderazgo en el proceso de fabricación, no ha logrado ningún éxito en chips de IA, su director ejecutivo fue destituido y sigue perdiendo participación de mercado en PC y servidores. El rendimiento de fabricación de Samsung también está en declive.

El escenario más extremo: el mundo podría terminar dependiendo únicamente de TSMC para la investigación y el desarrollo de chips de vanguardia. Esto significa que una isla, situada a unos cientos de kilómetros de la costa de China continental, carga con el destino de la tecnología global.

\subsection{Resumen de la sección}

La infraestructura de IA está atravesando una expansión sin precedentes: proyectos del calibre de Stargate representan inversiones del orden de \$100B. La electricidad, el enfriamiento y el contrabando de GPU son las nuevas restricciones físicas. El monopolio de NVIDIA no podrá revertirse en el corto plazo: el software de AMD es demasiado deficiente, Intel está en declive y el TPU de Google no está disponible externamente.
\section{La carrera de la IA, el código abierto y el panorama de la industria}

\subsection{¿Quién está ganando dinero?}

La única empresa que realmente está ganando dinero es NVIDIA. Las hiperescaladoras muestran utilidades contables, pero su gasto de capital es enorme. OpenAI va a la cabeza en el mejor modelo y en los ingresos de IA, pero gasta más de lo que ingresa. Meta se beneficia de la IA a través de sus sistemas de recomendación (no directamente de Llama).

\begin{warningbox}{La velocidad de la comoditización de las capacidades}
\textit{«Cualquier empresa cuyo modelo de negocio se basara en capacidades del nivel de GPT-3 está muerta. Cualquiera que se basara en el nivel de GPT-4 también está muerta.»} (Dylan)

OpenAI y Anthropic tienen que seguir ganando en el mejor modelo, o serán sustituidos por Llama o por modelos de código abierto. La curva de caída de costos es asombrosa: la inferencia de GPT-3 pasó de \$60-70 por millón de tokens a unos cuantos centavos de dólar en la actualidad, una caída de 1,200 veces en 3 años. DeepSeek está sobre la misma línea de tendencia de precios de los modelos del nivel de GPT-4; solo fue el primero en llegar.

Esto significa que el modelo de negocio de «modelo como servicio» enfrenta una presión de fondo. El único foso competitivo es mantenerse por delante de la frontera de manera continua: en cuanto te atrasas medio año, un modelo de código abierto alcanza tus capacidades y las vuelve gratuitas.
\end{warningbox}

\subsection{Agent: exageración vs. realidad}

«Agent» es el término de IA más abusado de 2025; su significado real debería ser la resolución de tareas abierta e independiente.

\begin{knowledgebox}{El problema de los seis nueves}
Cada paso tiene una precisión $<$100\%, y al componer varios pasos el rendimiento final cae mucho. Si cada paso tiene 99\% de precisión, después de 10 pasos solo queda 90\%; después de 100 pasos solo queda 37\%. ¿Cuántos nueves necesita cada paso para alcanzar una tasa de éxito de extremo a extremo de 99.99\%?

\textit{«¿Cuántos nueves tienes? Multiplícalos por el número de pasos... ni siquiera 99.9999\% alcanza.»} (Dylan)

La analogía con la conducción autónoma: incluso en carreteras bien definidas sigue siendo enormemente difícil (a Waymo le tomó diez años); un entorno abierto de páginas web o sistemas operativos es varias veces más caótico. Cada sitio web tiene un diseño distinto, las API cambian en cualquier momento y los patrones de error son innumerables.
\end{knowledgebox}

\begin{figure}[H]
\centering
\includegraphics[width=0.92\textwidth]{frame-05-agents.jpg}
\caption{Discusión sobre los límites de viabilidad y confiabilidad de los Agents\protect\footnotemark}
\end{figure}
\footnotetext{Intervalo de tiempo en el video: 04:17:00--04:17:12.}

La dirección más prometedora es el \textbf{Agent de ingeniería de software}, por las siguientes razones:
\begin{itemize}
    \item Verificabilidad: las pruebas unitarias, la compilación y el CI/CD ofrecen retroalimentación automática
    \item Puede revisar toda la base de código (la ventana de contexto es suficiente)
    \item SWE-bench pasó de 4\% a 60\% en un año
\end{itemize}

La colaboración estructurada también funciona (OpenAI Operator + DoorDash/OpenTable): en dominios bien definidos, la tasa de éxito de los Agent es mayor.

El costo de la ingeniería de software va a desplomarse, lo cual cambia toda la estructura del mercado. China no usa SaaS porque sus ingenieros son baratos, y la IA llevará ese modelo a todo el mundo. Los expertos de dominio (aviación, semiconductores, industria química) usan herramientas de hace 20 años: \textit{«Las herramientas de litografía de ASML corren en Windows XP.»} (Dylan) Estos ámbitos tienen mucha fruta al alcance de la mano.

\subsubsection{DOGE y la modernización del gobierno}

El software del gobierno está sumamente atrasado: \textit{«suplica una modernización»}. La burocracia protege centros de poder; el software rompe esas barreras. Muchas industrias tienen abundante fruta al alcance de la mano para la automatización con IA.

\subsection{La controversia de la destilación}

\begin{knowledgebox}{La zona gris ética de la destilación}
La destilación (entrenar un modelo más débil con las salidas de uno más fuerte) es una práctica estándar en la industria. OpenAI afirma que DeepSeek usó las salidas de su API, pero esos datos ya estaban en internet por medio de copiar y pegar. Muchas empresas emergentes estadounidenses entrenan abiertamente con salidas de OpenAI.

\textit{«¿Por qué es poco ético que yo entrene con las salidas de tu modelo, mientras que tú puedes entrenar con el texto de internet?»} (Nathan)

El espionaje industrial también existe: las ideas fluyen libremente a través de los cambios de empleo en Silicon Valley (en California las cláusulas de no competencia son ilegales). Un ingeniero de Gemini que ayudó a construir el context de 1 millón se fue después a Meta; se espera que la siguiente generación de Llama tenga un context de 1 millón.

El resquicio de Japón: la ley de derechos de autor permite entrenar con cualquier dato, más 9 GW de energía nuclear detenida, más importación ilimitada de GPU, es igual a un posible paraíso para el entrenamiento de IA.
\end{knowledgebox}

\subsection{El impacto de DeepSeek R1 en el código abierto}

La licencia MIT de R1 es un «gran reinicio»: el primer modelo de frontera con una licencia verdaderamente permisiva. Antes, las opciones eran modelos que no eran de frontera o licencias restrictivas (como la de Llama).

\textit{«Necesitamos modelos verdaderamente abiertos... no sabemos cuáles son los datos de DeepSeek R1, pero puedes usarlo para hacer una copia barata y hacer como si fuera tuya. Con Llama no puedes hacer eso.»} (Nathan)

El proyecto Tulu de AI2 de Nathan demuestra el poder de un posentrenamiento completamente abierto: al aplicar RLVR (RL con recompensas verificables) sobre la base de Llama, superó a Llama instruct en matemáticas e igualó a DeepSeek V3. Entre los modelos ubicados entre los primeros 60 lugares de Chatbot Arena, antes de Tulu \textbf{ninguno} había publicado el código o los datos de su posentrenamiento.

\subsection{Resumen de la sección}

El panorama de la industria de la IA se está consolidando con rapidez: solo NVIDIA es rentable, y la comoditización de las capacidades avanza a toda velocidad. Los Agent enfrentan el reto de confiabilidad de los seis nueves. El código abierto está cambiando las reglas del juego: la licencia MIT de DeepSeek R1 y el proyecto Tulu de AI2 demuestran que un posentrenamiento completamente abierto puede alcanzar el nivel de frontera. La controversia de la destilación refleja un problema sin resolver sobre la propiedad de los datos.
\section{Caída de costos, la paradoja de Jevons y el contrabando industrial}

\subsection{La caída de costos de 1200 veces y la paradoja de Jevons}

El costo de inferencia en la era de GPT-3 era de \$60-70 por millón de tokens. Tres años después (mediante modelos pequeños como Llama 3B), la misma capacidad bajó a unos cuantos centavos: una reducción de \textbf{1200 veces}. DeepSeek está en la línea de tendencia de costos de un modelo del nivel de GPT-4, pero fue el primero en llegar a ese punto de precio.

Las acciones de NVIDIA se desplomaron tras el lanzamiento de DeepSeek, pero eso obedeció a un contagio social basado en una narrativa equivocada. Ningún modelo público ha costado más de \$1000 millones en entrenamiento. La afirmación de que «DeepSeek solo costó 6 millones de dólares» es una mala lectura del artículo: eso es únicamente el costo de horas de GPU de una sola corrida de preentrenamiento, sin incluir investigación y desarrollo, experimentación, entrenamiento posterior ni despliegue de inferencia.

\begin{importantbox}{La paradoja de Jevons confirmada a la perfección en la IA}
La paradoja de Jevons en economía (planteada en 1865): cuando mejora la eficiencia con la que se usa un recurso, el consumo total aumenta en lugar de disminuir, porque el menor costo estimula más demanda.

La eficiencia de la máquina de vapor mejora $\to$ el consumo de carbón no baja, sube. El costo de inferencia de IA baja $\to$ la demanda de GPU no baja, sube.

Evidencia directa tras el lanzamiento de DeepSeek:
\begin{itemize}
    \item El precio de renta de H100 en AWS no bajó, \textbf{subió}
    \item El H200 casi se agotó
    \item Todos los proveedores de nube reportan una demanda disparada de GPU
    \item Nuevos casos de uso (agentes intensivos en inferencia, RAG a gran escala) se volvieron viables económicamente
\end{itemize}

\textit{``'Las Scaling Laws murieron' duró un mes... y luego salieron O1, O3, R1, y ahora se volvió 'los modelos avanzan demasiado rápido'.''} (Dylan)

Inferencia más barata $\to$ más casos de uso $\to$ más demanda de inferencia $\to$ se necesitan más GPU. No es un juego de suma cero: el mercado total se está expandiendo. Por eso las perspectivas de largo plazo de NVIDIA no se oscurecen por DeepSeek.
\end{importantbox}

\subsection{La escala y los métodos del contrabando de GPU}

El contrabando de GPU pasó de ser una zona gris a una operación de nivel industrial, con múltiples canales:

\textbf{Canal de renta en la nube}: ByteDance es el «contrabandista» más grande: renta más de 500,000 GPU a nivel mundial a través de Oracle, Google y diversos proveedores de nube pequeños. Reparte la operación entre entidades en varios países para evitar el escrutinio de una sola jurisdicción. Este «contrabando distribuido» es extremadamente difícil de rastrear porque cada contrato de renta individual queda por debajo del umbral regulatorio.

\textbf{Contrabando físico}: alguien viajó en primera clase desde el aeropuerto de San Francisco a Shanghái con cajas de servidores GPU de SuperMicro; esto no es ficción, ocurrió de verdad. Se estima que en 2024 entre 200,000 y 300,000 GPU llegaron a China a través de pequeñas empresas en Singapur y Malasia.

Los límites de las \textbf{nuevas reglas de difusión}: restringen a las entidades vinculadas con China a rentar menos de 2000 GPU o comprar menos de 1500 GPU. Pero el reto de aplicación es enorme: ¿cómo rastrear al usuario final de recursos de cómputo en la nube a escala global? ¿Cómo se define una «entidad vinculada con China»? ¿Una empresa registrada en Singapur pero operada por ciudadanos chinos cuenta o no?

\textit{``Las GPU van a superar a las drogas y a las armas como el bien de contrabando de mayor valor por kilogramo.''} (Dylan) Una H100 cuesta cerca de \$30,000 y pesa menos de 5 kg, más de \$6,000 por kilogramo, muy por encima de la mayoría de los bienes prohibidos. Y a diferencia de las drogas, las GPU son mercancía legal en algunos mercados, solo restringida en ciertas direcciones de exportación. Eso crea un espacio de arbitraje singular.

\subsection{Destilación, espionaje y flujo de conocimiento}

\subsubsection{La naturaleza de la controversia sobre la destilación}

La destilación (entrenar un modelo más débil con las salidas de uno más fuerte) es una práctica estándar en la industria. OpenAI afirmó que DeepSeek usó las salidas de su API, y el Financial Times lo publicó en un titular. Pero Nathan señala que esto se parece más a un «control de la narrativa» que a una acusación técnica real:

\textit{``¿Por qué es poco ético que yo entrene con las salidas de tu modelo, mientras que tú puedes entrenar con el texto de internet?''} (Nathan) El contenido generado por usuarios a través de ChatGPT se copia y pega en internet, y así se vuelve un dato público. Muchas empresas emergentes estadounidenses entrenan abiertamente sobre salidas de OpenAI; ¿por qué ellas sí pueden y DeepSeek no?

Una pregunta más profunda: \textbf{¿a quién pertenece la propiedad intelectual de las salidas de un modelo?} Si yo genero un fragmento de código con GPT-4, ¿el derecho de autor de ese código es de OpenAI o mío? Esta pregunta todavía no tiene una respuesta definitiva en el terreno legal.

\subsubsection{Movilidad de talento y espionaje industrial}

La cultura de cambio de empleo de Silicon Valley (las cláusulas de no competencia son ilegales en California) hace que las ideas —no el código— circulen libremente. Un caso concreto: un ingeniero que ayudó a Google a construir el contexto de un millón de tokens de Gemini se cambió a Meta, y se espera que la siguiente generación de Llama tenga un contexto de un millón de tokens. ¿Esto es transferencia legítima de conocimiento o robo de tecnología de facto? El límite es sumamente difuso.

Robar código o datos es un delito y es relativamente fácil de rastrear. Pero \textbf{la transferencia de ideas} es inevitable: no se puede «olvidar» la intuición arquitectónica que se aprendió en la empresa anterior.

\begin{knowledgebox}{Japón: un paraíso inesperado para el entrenamiento de IA}
Japón podría convertirse en un nodo singular para el entrenamiento de IA a nivel mundial:
\begin{itemize}
    \item \textbf{Ley de derechos de autor}: permite explícitamente entrenar modelos con cualquier dato (una exención legal poco común)
    \item \textbf{Energía}: 9 GW de capacidad nuclear inactiva que puede reactivarse para centros de datos, una enorme fuente de electricidad barata
    \item \textbf{Hardware}: sin restricciones de importación de GPU (no está sujeto a los controles de exportación de Estados Unidos porque Japón es un país aliado)
    \item Varias empresas japonesas y capital extranjero ya están construyendo grandes instalaciones de entrenamiento de IA en Japón
\end{itemize}

Exención de derechos de autor + energía nuclear barata + GPU sin límite = un entorno de entrenamiento ideal. Es una oportunidad de arbitraje geográfico que casi nadie ha notado.
\end{knowledgebox}

También ocurren actividades de inteligencia dirigidas a investigadores de IA. Dylan señala, medio en broma: \textit{``como hombre soltero de veintitantos años... somos muy fáciles de corromper.''} Las operaciones de trampa (honey pot) dirigidas a jóvenes investigadores con conocimiento de frontera son una amenaza de seguridad real. Los servicios de inteligencia de distintos países compiten por el talento en IA, no para reclutar espías, sino para obtener secretos de entrenamiento y conocimientos arquitectónicos.

\subsection{La misión de código abierto de Nathan en AI2/Tulu}

El trabajo de Nathan en AI2 (Allen Institute for AI) representa el estándar más alto de IA de código abierto. Su proyecto Tulu fue el primero, entre los modelos que están en los 60 primeros lugares de Chatbot Arena, en publicar la receta \textbf{completa} de entrenamiento posterior (código + datos).

Resultado clave: Tulu aplica RLVR (RL con recompensas verificables) sobre el modelo base Llama, supera a la versión instruct oficial de Llama en razonamiento matemático y iguala el desempeño de DeepSeek V3. Esto demuestra que el \textbf{método de entrenamiento posterior en sí mismo} (y no solo la escala del modelo base) puede producir una diferencia enorme de desempeño.

OLMo es otro proyecto de AI2: un trabajo de preentrenamiento totalmente abierto (que complementa el entrenamiento posterior de Tulu). Con una escala de modelo mayor, es más fácil que el entrenamiento con RL despierte capacidades sólidas, que luego se pueden destilar a modelos pequeños.

La motivación central de Nathan: \textit{``No confío en la gente que dice 'confía en mí, hermano, vamos a hacer que la IA sea buena'.''} La IA es la tecnología más poderosa de nuestra vida: se necesita que más personas participen en darle forma, en lugar de confiar en que unas cuantas empresas decidan a puerta cerrada. La apertura ayuda a que investigadores fuera del campo de la IA, gobiernos y todos en general puedan entender lo que está ocurriendo.

\subsection{Resumen de la sección}

La economía de la industria de la IA se está reconfigurando a gran velocidad: la caída de costos de 1200 veces queda contrarrestada por la paradoja de Jevons, y la demanda total sigue creciendo. El contrabando de GPU se convirtió en una nueva industria subterránea. La destilación y la movilidad de talento hacen que los límites de la propiedad intelectual sean extremadamente difusos. La combinación singular de Japón de ley, energía y hardware crea una ventaja inesperada para el entrenamiento. El trabajo de Nathan en AI2/Tulu demuestra que el entrenamiento posterior totalmente abierto puede alcanzar el nivel de frontera.
\section{Métodos de evaluación: de la puntuación del modelo a la capacidad del sistema}

\subsection{Por qué un benchmark puntual suele ser engañoso}

En la entrevista aparece una y otra vez la misma señal: muchos debates que parecen tratar sobre si «el modelo es más fuerte» en realidad tratan sobre si «el sistema se puede implementar en producción». Estas dos cosas ya se separaron claramente en 2025. Que un modelo esté a la cabeza en un benchmark matemático no significa que pueda sustituir a un humano en un flujo de trabajo real; que un modelo tenga una clasificación alta en Chatbot Arena tampoco significa que pueda completar de manera estable tareas que abarcan varios sistemas.

La prueba subjetiva que Lex presenta en el programa es representativa: en preguntas de filosofía, O1 Pro queda establemente en primer lugar, R1 en segundo y Gemini Flash Thinking en tercero. Pero eso es solo la dimensión de «respuesta de alta calidad en una sola ronda». Las dimensiones que más les interesan a Dylan y Nathan son:
\begin{itemize}
    \item Costo por tarea (no solo el precio unitario por token, sino también los reintentos y la supervisión)
    \item Tasa de éxito en tareas largas (tasa de finalización de extremo a extremo en procesos de 10 a 100 pasos)
    \item Capacidad de diagnóstico de los modos de falla (si, tras una falla, se puede ubicar con rapidez si el problema es del Prompt, del Tool, del modelo o del sistema)
    \item Capacidad de rendimiento bajo cómputo limitado (concurrencia, latencia, ocupación de KV Cache)
\end{itemize}

\begin{knowledgebox}{«La evaluación de modelos» y «la evaluación de sistemas» deben separarse por capas}
\textbf{Capa de modelo (Model-level)}: MMLU, AIME, benchmark de código, puntaje de preferencia en Arena. Sirve para comparar el techo de capacidad.\\
\textbf{Capa de sistema (System-level)}: tasa de éxito de la tarea de extremo a extremo, número promedio de reintentos, latencia P95, costo por resultado. Sirve para juzgar la viabilidad del negocio.\\
\textbf{Capa organizacional (Org-level)}: velocidad de corrección, eficiencia del retorno de datos, cadencia de lanzamientos, estabilidad frente a regresiones. Sirve para juzgar si el liderazgo es sostenible.

El planteamiento central de la entrevista no es «quién queda primero en la tabla», sino «quién puede convertir la capacidad en producción estable». Por eso hacen énfasis en la infraestructura, el suministro eléctrico y los procesos de ingeniería, y no solo en el número de parámetros.
\end{knowledgebox}

\subsection{Evaluación consciente del costo: poner la exactitud y los dólares en una misma tabla}

Muchos equipos, en la etapa de PoC, solo observan la exactitud e ignoran el cómputo de inferencia y el costo de los reintentos, y solo al llegar a producción descubren que la economía unitaria no es sostenible. Combinando esto con la discusión del programa, un marco de evaluación más práctico es el siguiente:

\begin{table}[H]
\centering
\begin{tabular}{lp{9cm}}
\toprule
\textbf{Indicador} & \textbf{Definición y recomendación práctica} \\
\midrule
Tasa de aprobación en un intento (Pass@1) & Proporción de tareas completadas directamente en una sola ejecución; refleja el techo de la experiencia real \\
Tasa de aprobación tras reintentos (Pass@k) & Tasa de finalización tras permitir $k$ reintentos; refleja la capacidad de «cambiar dinero por tasa de éxito» \\
Costo por éxito & \texttt{costo por llamada $\times$ número promedio de llamadas / número de tareas exitosas} \\
Latencia por éxito & Tiempo total desde el inicio de la tarea hasta obtener un resultado utilizable (incluidos los reintentos) \\
Proporción de respaldo humano & Proporción de casos que requieren que un humano tome el control; determina directamente la escalabilidad de la organización \\
\bottomrule
\end{tabular}
\caption{Indicadores de evaluación consciente del costo orientados a la implementación}
\end{table}

Si solo se observa Pass@1, los modelos de razonamiento pueden parecer «más lentos y más caros»; pero si se observan Pass@k y la proporción de respaldo humano, los modelos de razonamiento pueden resultar más económicos en el costo total para tareas complejas. El valor de modelos como R1 y O1 por lo general no está en «una llamada barata», sino en «reducir las correcciones repetidas por parte de una persona».

\begin{importantbox}{Las conclusiones de la evaluación deben escribirse como una tripleta «capacidad-costo-riesgo»}
Una plantilla de conclusión ejecutable es la siguiente:
\begin{itemize}
    \item \textbf{Capacidad}: en qué tareas supera de manera significativa a la línea base (por ejemplo, corrección de código en varios pasos, demostraciones matemáticas)
    \item \textbf{Costo}: cuánto presupuesto y latencia adicionales cuesta cada punto porcentual de mejora en la tasa de éxito
    \item \textbf{Riesgo}: cuál es el modo de falla más común, cuáles son sus condiciones de activación y si existe un mecanismo de respaldo
\end{itemize}

Solo cuando los tres elementos se escriben juntos, la evaluación del modelo puede orientar las decisiones de lanzamiento a producción, en lugar de ser solo un reporte de benchmark bien presentado.
\end{importantbox}

\subsection{Confiabilidad en tareas largas: cuántos «nueves» bastan}

El «how many nines» de la entrevista es, en esencia, una multiplicación de confiabilidad. Si la tasa de éxito de cada paso es $p$ y la tarea tiene en total $n$ pasos, la tasa de éxito de extremo a extremo es aproximadamente $p^n$. Esto explica por qué un Agent tiene un desempeño aceptable en cadenas cortas, pero falla rápidamente en cadenas largas.

\begin{table}[H]
\centering
\begin{tabular}{lccc}
\toprule
\textbf{Tasa de éxito por paso} & \textbf{Tarea de 10 pasos} & \textbf{Tarea de 30 pasos} & \textbf{Tarea de 100 pasos} \\
\midrule
99\% & 90.4\% & 73.9\% & 36.6\% \\
99.5\% & 95.1\% & 86.1\% & 60.6\% \\
99.9\% & 99.0\% & 97.0\% & 90.5\% \\
\bottomrule
\end{tabular}
\caption{Efecto compuesto de la tasa de éxito en tareas de varios pasos}
\end{table}

Por eso Dylan insiste una y otra vez: lo importante en un Agent de entorno abierto no es «si sabe dar el siguiente paso», sino «si sigue siendo controlable después de decenas de pasos». Por lo tanto, la ruta realmente práctica suele consistir primero en dividir la tarea en subtareas verificables, y después usar un orchestrator para controlar los límites de entrada y salida de cada paso.

\begin{warningbox}{Evitar confundir un «Demo exitoso» con un «sistema confiable»}
En la entrevista se ven muchos casos exitosos, pero los errores de juicio comunes al llevarlo a la ingeniería en producción son:
\begin{itemize}
    \item Ajustar el prompt repetidamente sobre el mismo conjunto de datos y, al obtener una puntuación alta, lanzarlo directamente a producción
    \item Ignorar el aumento súbito de fallas causado por la deriva del entorno (rediseños de páginas web, actualizaciones de API, cambios de permisos)
    \item Contabilizar solo los casos exitosos, sin registrar las trayectorias de falla ni el costo de recuperación
\end{itemize}

Para tareas de varias etapas que duran más de una hora, se recomienda diseñar la «reversión automática tras una falla» y los «puntos de relevo humano» como elementos de primera clase, y no como parches posteriores.
\end{warningbox}

\subsection{Resumen de la sección}

La evaluación no debe quedarse en el nivel de la puntuación, sino avanzar hacia una medición conjunta de «capacidad-costo-riesgo». Para los modelos de razonamiento y los sistemas de Agent, la confiabilidad de extremo a extremo y el costo por éxito importan más que un benchmark puntual. La distorsión multiplicativa de las tareas largas determina que, sin mecanismos de verificación ni procesos de respaldo, no hay implementación sostenible.
\section{Organización y ejecución: por qué DeepSeek generó una ventaja de velocidad}

\subsection{El ciclo integrado de investigación, sistemas e infraestructura}

La entrevista da, de manera reiterada, una conclusión implícita: la ventaja de DeepSeek no es un punto algorítmico aislado, sino la capacidad de coordinación entre capas. MoE, MLA, la optimización de la comunicación de bajo nivel y la estrategia de servicio de inferencia no son el resultado de equipos aislados que trabajan por separado y luego se ensamblan, sino que convergen en torno al mismo cuello de botella (la potencia de cómputo limitada).

\begin{table}[H]
\centering
\begin{tabular}{lp{8.5cm}}
\toprule
\textbf{Nivel} & \textbf{Enfoque al estilo DeepSeek (según la entrevista)} \\
\midrule
Nivel del modelo & Usar MoE para reducir los parámetros activados y MLA para aliviar la presión de la KV Cache \\
Nivel del sistema & Optimizar la ruta de comunicación por debajo de CUDA/NCCL, para reducir las pérdidas causadas por las limitaciones de interconexión \\
Nivel del servicio & Priorizar, del lado de la inferencia, la ampliación del beneficio en tareas verificables, sacrificando parte de la experiencia general a cambio de una mejor relación costo-beneficio \\
Nivel organizacional & Permitir que investigación e ingeniería iteren con rapidez en torno al mismo objetivo, en lugar de entregar de manera lineal \\
\bottomrule
\end{tabular}
\caption{La ventaja de DeepSeek se parece más a una «coordinación entre capas» que a una innovación puntual}
\end{table}

Este tipo de coordinación es difícil de replicar mediante la simple compra de recursos. Comprar las mismas GPU no equivale a producir el mismo rendimiento; la diferencia real está en si el equipo puede convertir la «innovación de artículo» en una «mejora del rendimiento de extremo a extremo».

\subsection{De Ablation a YOLO Run: la ingeniería de las decisiones de alto riesgo}

El «YOLO run» que describe Nathan es el paradigma típico del entrenamiento de frontera: se valida la dirección con ablaciones a pequeña escala y, después, en una ventana de tiempo determinada, se concentran los recursos en una sola ejecución de entrenamiento a gran escala. Suena a una apuesta, pero en realidad su condición previa es una disciplina experimental estricta y la capacidad de hacer una revisión posterior con rapidez.

\begin{knowledgebox}{Un manual operativo aplicable para el «YOLO run»}
\begin{itemize}
    \item Definir con claridad el único objetivo de optimización de esta ronda (por ejemplo, verificar si una arquitectura determinada aumenta los tokens efectivos dentro de un presupuesto de cómputo fijo)
    \item Definir de antemano las condiciones de interrupción (anomalías en la loss, explosión del gradiente, desviación del rendimiento respecto al umbral)
    \item Preparar rutas de reversión para las fallas comunes (estrategia de checkpoint, cambio de datos, degradación de la precisión mixta)
    \item Completar el postmortem dentro de las 24 a 48 horas posteriores al fin del entrenamiento, y actualizar la lista de ablation de la siguiente ronda
\end{itemize}

Sin estos mecanismos, el «YOLO» es solo una apuesta arriesgada; con ellos, el «YOLO» se convierte en una forma de investigación y desarrollo de alto apalancamiento que se puede gestionar.
\end{knowledgebox}

En el programa, ambos mencionan que, durante el entrenamiento, vigilan de manera continua la curva de loss y atienden repetidamente los spikes. Estos detalles muestran que la capacidad de frontera no depende de un golpe de inspiración, sino de un proceso de ingeniería de alta intensidad, reproducible y corregible.

\subsection{Estrategia de lanzamiento: la velocidad misma es un arma competitiva}

«DeepSeek ships» es una frase clave en la entrevista. El lanzamiento rápido no es una consigna de mercadotecnia, sino un mecanismo de aprendizaje: entrar antes al tráfico real $\to$ descubrir antes los modos de falla $\to$ corregir con mayor rapidez $\to$ entrar con mayor rapidez a la siguiente ronda de optimización.

Esta estrategia tiene un costo: riesgo de marca, controversias de seguridad, presión operativa. La disyuntiva de DeepSeek consiste en sacrificar parte de la estabilidad a cambio de velocidad de aprendizaje, mientras que la disyuntiva de Anthropic es casi la opuesta. Ambas rutas son coherentes en sí mismas, pero corresponden a objetivos organizacionales distintos.

\begin{warningbox}{La ventaja de velocidad es lo más fácil de destruir con KPI equivocados}
Muchas organizaciones persiguen la velocidad de palabra, pero en la práctica usan «cero incidentes» y «cada lanzamiento debe ser perfecto» como KPI centrales, y con eso terminan por bloquear la velocidad de iteración. La realidad que muestra la entrevista es que, en un periodo de competencia de frontera, «falla controlada + recuperación rápida» suele tener más probabilidad de éxito que «buscar acertar a la primera».
\end{warningbox}

\subsection{Resumen de la sección}

La ventaja de velocidad de DeepSeek proviene de la coordinación entre capas, de la ingeniería de las decisiones de alto riesgo y del ciclo de retroalimentación del lanzamiento. No es una técnica que se pueda copiar en un solo punto, sino el resultado conjunto de la estructura organizacional, la disciplina de investigación y desarrollo, y la asignación de recursos.
\section{Lista de implementación: empresas, equipos de investigación y quienes formulan políticas}

\subsection{Para quien aplica esto en la empresa: reforzar primero el «proceso verificable»}

Para la mayoría de las empresas, el orden correcto no es «perseguir primero el modelo más reciente», sino «definir primero una tarea verificable». En la entrevista se menciona varias veces que el código y las matemáticas avanzan rápido porque cuentan con un verificador de aceptación claro (tests / checker).

\begin{importantbox}{Ciclo mínimo de 6 pasos para la implementación empresarial}
\begin{enumerate}
    \item Elegir 1 o 2 procesos de alto valor con aceptación automatizable (por ejemplo, corrección de código, extracción de documentos, enrutamiento de tickets)
    \item Establecer un conjunto de evaluación unificado, con un modelo base y un Prompt base fijos
    \item Adoptar Pass@k más el costo por éxito como métricas centrales
    \item Registrar de manera obligatoria las trayectorias de fallo (entrada, llamadas a herramientas, tipo de error, ruta de recuperación)
    \item Diseñar puntos de intervención humana y un mecanismo de reversión
    \item Revisar semanalmente el Top 10 de fallos antes de decidir si se cambia de modelo
\end{enumerate}

Si desde el inicio se persigue un «Agente totalmente autónomo», por lo general se cae en una doble trampa de confiabilidad y gobernanza.
\end{importantbox}

\subsection{Para los equipos de modelos fundacionales: convertir la «ventaja de entrenamiento» en «ventaja de servicio»}

La entrevista señala una fractura frecuente: el modelo va a la delantera en artículos y benchmarks, pero la concurrencia, la latencia y la disponibilidad del lado del servicio no la siguen, y al final el valor lo capta una plataforma de alojamiento externa. Para evitar esta fractura, hay que optimizar al menos tres cosas de manera simultánea:
\begin{itemize}
    \item \textbf{Arquitectura de inferencia}: estrategia de KV Cache, estrategia de procesamiento por lotes, programación de contexto largo
    \item \textbf{Contrato de producto}: qué capacidades se pueden prometer y cuáles solo se abren de forma experimental
    \item \textbf{Experiencia del desarrollador}: interpretabilidad de errores, transparencia de cuotas, ritmo de las actualizaciones de versión
\end{itemize}

\begin{table}[H]
\centering
\begin{tabular}{lp{8.5cm}}
\toprule
\textbf{Objetivo} & \textbf{Acción recomendada} \\
\midrule
Reducir el costo de las alucinaciones & Exigir el uso de herramientas más un checker en las rutas verificables, y reducir los ciclos cerrados que dependen solo del lenguaje natural \\
Reducir el costo de los reintentos & Emitir códigos de error estructurados y las causas del fallo, para reducir los reintentos a ciegas \\
Elevar la eficiencia de la migración & Al lanzar un nuevo modelo, ofrecer una capa de compatibilidad y un interruptor de reversión, para que el desarrollador no tenga que reescribir todo de una vez \\
\bottomrule
\end{tabular}
\caption{Acciones clave para pasar de «liderazgo del modelo» a «liderazgo de la plataforma»}
\end{table}

\subsection{Para quienes formulan políticas: convertir los objetivos regulatorios en indicadores medibles}

La lección más importante del programa sobre los controles de exportación es esta: si el objetivo se limita a «restringir los flops de entrenamiento», en el nivel de ejecución surgirán numerosas rutas alternativas (memoria, interconexión, arrendamiento, entidades desdobladas). La política necesita pasar de la «restricción nominal» al «resultado medible».

Un marco de política más ejecutable podría incluir:
\begin{itemize}
    \item Indicadores de capacidad de entrenamiento: umbral de acceso a cómputo, disponibilidad de componentes clave
    \item Indicadores de capacidad de despliegue: capacidad de operar de manera sostenida clústeres de inferencia a gran escala
    \item Indicadores de capacidad de difusión: transparencia de las cadenas de arrendamiento y reventa transfronterizas
    \item Indicadores de capacidad de riesgo: resiliencia de la infraestructura crítica (red eléctrica, telecomunicaciones) frente a ataques con IA
\end{itemize}

\begin{warningbox}{No basta con vigilar «si entraron o no los chips»; también hay que vigilar «si se formó o no la capacidad»}
Si la evaluación de la política solo observa el flujo del hardware, puede resultar estadísticamente «en cumplimiento» y, en términos de capacidad, «desprotegida». Los casos reales que aporta la entrevista (arrendamiento distribuido, transbordo en zonas grises) muestran que los indicadores de la política deben vincularse con la capacidad finalmente formada; de lo contrario, es fácil caer en una falsa tranquilidad.
\end{warningbox}

\subsection{Resumen de la sección}

Ya se trate de empresas, equipos de modelos o quienes formulan políticas, lo importante es convertir el objetivo en un ciclo cerrado verificable, iterable y auditable. La competencia en IA no es un enfrentamiento de una sola vez, sino un proceso continuo de convertir la ventaja técnica en una ventaja de ejecución estable.
\section{Proyección de escenarios para los próximos 24 meses}

\subsection{Escenario A: el modelo avanza con rapidez, pero el despliegue sigue limitado por restricciones físicas}

Este es el juicio más consistente en la entrevista: la capacidad seguirá aumentando, pero el ritmo de despliegue está limitado por la energía eléctrica, el enfriamiento, la cadena de suministro y la gobernanza organizacional, por lo que no se puede «reescribir de la noche a la mañana todas las industrias». En este escenario, los primeros beneficiados no son las aplicaciones más futuristas, sino los flujos de trabajo verificables, integrables y auditables.

\begin{table}[H]
\centering
\begin{tabular}{lp{8.8cm}}
\toprule
\textbf{Capa de la industria} & \textbf{Cambio más probable en ocurrir primero} \\
\midrule
Desarrollo de software & Aumento continuo de la tasa de automatización en la corrección de código, la generación de pruebas y la refactorización de migraciones \\
Servicios empresariales & La clasificación de tickets, la extracción de documentos y la revisión de cumplimiento entran en una normalidad semiautomatizada \\
Manufactura y cadena de suministro & La simulación de planeación, la localización de anomalías y el ajuste de parámetros se realizan mediante colaboración entre humanos y máquinas \\
Sector público & La modernización de sistemas históricos avanza con lentitud, pero surgen rupturas puntuales en escenarios transaccionales de alta frecuencia \\
\bottomrule
\end{tabular}
\caption{Prioridades de evolución industrial en el escenario A}
\end{table}

En este escenario, «quién accede primero al modelo más potente» no es la ventaja decisiva; «quién logra integrar el modelo a un proceso estable» es la fuente real de la ventaja. El foso competitivo de las empresas se traslada de la «capacidad de adquisición de modelos» a la «capacidad de reestructuración de procesos».

\subsection{Escenario B: la fragmentación geopolítica se acelera y surge un ecosistema tecnológico de doble vía}

La discusión en la entrevista sobre los controles de exportación y las reglas de difusión implica un riesgo de mediano plazo: el ecosistema global de IA podría bifurcarse en una doble vía, «parcialmente interoperable, parcialmente no interoperable». El hardware, los recursos de nube, las licencias de modelos y el cumplimiento de datos formarían bloques regulatorios distintos.

\begin{warningbox}{Costos ocultos en un ecosistema de doble vía}
Las empresas suelen subestimar tres tipos de costos:
\begin{itemize}
    \item \textbf{Costo de cambio de cumplimiento}: la misma capacidad debe mantener distintas cadenas de suministro y procesos de auditoría según la región
    \item \textbf{Costo de cambio de modelo}: las licencias y los contratos de API son inconsistentes, lo que dispara el costo de migración
    \item \textbf{Costo de coordinación de talento}: a los equipos de distintas regiones les resulta difícil compartir el mismo conjunto de herramientas y el mismo circuito de datos
\end{itemize}

Esto significa que el «respaldo con múltiples modelos» y la «recuperación ante desastres multinube» no son opciones, sino infraestructura básica para la resiliencia operativa de mediano plazo.
\end{warningbox}

Si la bifurcación se acelera, la importancia de los modelos de código abierto aumentará todavía más, porque ofrecen una capa de amortiguación tanto en compatibilidad técnica como en poder de negociación. Una de las razones por las que Nathan enfatiza abrir las recetas de posentrenamiento es precisamente dotar al ecosistema de capacidad de resistencia frente al monopolio y al bloqueo.

\subsection{Escenario C: un incidente de seguridad detona un salto regulatorio}

Si ocurre un incidente de gran impacto (por ejemplo, un ataque a infraestructura crítica asistido por IA a gran escala), el ritmo regulatorio podría pasar de gradual a discontinuo. En ese momento el mercado experimentaría una «contracción de capacidad» de corto plazo, pero a mediano y largo plazo esto reforzaría el valor de los sistemas auditables.

\begin{importantbox}{Capacidades que deben prepararse antes de un salto regulatorio}
\begin{itemize}
    \item Registros de auditoría a nivel de tarea (entradas, versión del modelo, llamadas a herramientas, salidas)
    \item Cadena de evidencia explicable para las decisiones clave (por qué se sugirió esa acción)
    \item Mecanismo automático de corte por nivel de riesgo (las tareas de alto riesgo activan una aprobación manual)
    \item Capacidad de rastreo de versiones de modelos y datos (para ubicar el origen de un problema)
\end{itemize}

Los equipos que cuentan con estas capacidades, ante un cambio regulatorio, no quedan paralizados de forma pasiva, sino que pueden pasar la auditoría con mayor rapidez y reanudar operaciones.
\end{importantbox}

\subsection{Resumen de la sección}

Lo más probable en los próximos dos años es una tensión estructural de «la capacidad avanza rápido, el sistema se implementa despacio». La fragmentación geopolítica y los saltos regulatorios magnificarán la importancia de la capacidad de ejecución. La verdadera competencia no es solo la fortaleza del modelo, sino quién logra mantener una entrega continua en un entorno incierto.
\section{Plano de ingeniería: conectar los modelos de razonamiento a sistemas de producción}

\subsection{Arquitectura por capas: Planner / Executor / Verifier}

Combinando lo que se discutió en la entrevista sobre tareas verificables, confiabilidad en cadenas largas y restricciones de costo, un plano de ingeniería práctico es una arquitectura de tres capas:
\begin{itemize}
    \item \textbf{Planner}: se encarga de descomponer la tarea, elegir herramientas y controlar el presupuesto
    \item \textbf{Executor}: invoca el modelo y las herramientas para ejecutar las subtareas
    \item \textbf{Verifier}: aplica validación por reglas, validación mediante pruebas o validación de consistencia al producto de cada paso
\end{itemize}

\begin{knowledgebox}{Por qué la arquitectura de tres capas es más estable que un «modelo único que lo hace todo»}
Un modelo único de extremo a extremo se desarrolla rápido, pero es difícil de controlar en cuanto a los límites de la falla. Al introducir un Verifier, se puede contener el error en la capa de la subtarea, en lugar de dejar que se propague hasta el resultado final. En tareas de más de 30 pasos, esta diferencia se amplifica exponencialmente.
\end{knowledgebox}

En este plano, los modelos de razonamiento fuerte deben usarse de preferencia en los pasos de «alta incertidumbre y alto costo»; los modelos de bajo costo se encargan de los pasos de «alta frecuencia y bajo riesgo». Así se mantiene la calidad y a la vez se controla el presupuesto total.

\subsection{Contexto y memoria: evitar que «entre más se avanza, más se desordena»}

La causa común de que fallen las tareas largas no es que el modelo no sepa hacerlas, sino la contaminación del contexto: información errónea temprana que se toma como verdad y se sigue propagando. La discusión de la entrevista sobre la KV Cache y el costo del contexto largo muestra que un «contexto infinito» no resuelve automáticamente el problema de la calidad de la memoria.

\begin{table}[H]
\centering
\begin{tabular}{lp{8.7cm}}
\toprule
\textbf{Problema} & \textbf{Medida de ingeniería} \\
\midrule
Expansión del contexto & Usar resúmenes por etapas e instantáneas de estado, en lugar de acumular el historial sin límite \\
Consolidación de memoria errónea & Asignar niveles de confianza a los hechos clave; la información de baja confianza debe verificarse una segunda vez \\
Desplazamiento semántico entre herramientas & Usar representaciones intermedias estructuradas (JSON schema / estado tipado) \\
Contaminación por reintentos & Cada reintento debe reproducirse desde un estado limpio, evitando acumular parches sobre una trayectoria errónea \\
\bottomrule
\end{tabular}
\caption{Estrategias de gobernanza del contexto en tareas largas}
\end{table}

\subsection{Estrategia de despliegue: liberación gradual, reversión y división del trabajo entre humano y máquina}

El despliegue de los modelos de razonamiento no debe hacerse como «reemplazo total», sino como «liberación gradual por escenario». Se puede estratificar por nivel de riesgo:
\begin{itemize}
    \item L1 (riesgo bajo): puede ejecutarse de forma automática, con revisión humana por muestreo
    \item L2 (riesgo medio): el modelo propone y el humano confirma
    \item L3 (riesgo alto): el modelo solo entrega el análisis, sin ejecutar directamente
\end{itemize}

\begin{warningbox}{La forma más peligrosa de desplegar: tratar una tarea L3 como si fuera L1}
Cuando una organización, bajo presión de KPI, persigue prematuramente una alta tasa de automatización, el error más común es ejecutar de forma automática y directa procesos de alto riesgo. Una vez que este tipo de error se dispara, la consecuencia no suele ser una «caída de precisión», sino un resultado «a nivel de accidente». La preocupación que muestra la entrevista sobre los ataques cibernéticos y la vulnerabilidad del sistema eléctrico, en esencia, es un recordatorio de este límite.
\end{warningbox}

Además, la reversión debe poder «ensayarse», y no quedarse solo escrita en la documentación. Se recomienda hacer un simulacro de falla una vez al mes: inyectar de manera aleatoria salidas anómalas del modelo y verificar si el sistema puede volver a la ruta manual dentro de un tiempo limitado.

\subsection{Resumen de la sección}

Llevar los modelos de razonamiento a producción no es solo ingeniería de prompts, sino ingeniería de sistemas. La arquitectura por capas, la gobernanza del contexto, el despliegue estratificado por riesgo y la reversión ensayable determinan si la capacidad del modelo se convierte o no en productividad estable.

\section{Anécdotas, teorías de conspiración y jugadores subestimados}

Esta sección recopila ideas y anécdotas dispersas pero muy valiosas de la entrevista: no pertenecen a un solo tema, pero cada una revela algún aspecto de la industria de la IA.

\subsection{Teorías de conspiración contra la realidad}

Las teorías de conspiración en torno a DeepSeek se propagaron como fuego en las redes sociales. Dylan y Nathan las desmontaron una por una:

\textbf{«El gobierno chino está subsidiando a DeepSeek»}: poco probable. DeepSeek está financiada por High-Flyer Quant (幻方量化), un fondo de cobertura, no por una empresa vinculada al gobierno. Liang Wenfeng (梁文锋) sí se reunió después con dirigentes del gobierno, pero eso fue \textbf{después} del éxito de DeepSeek: el gobierno se sumó a la ola, no la financió desde atrás.

\textbf{«Vendieron en corto las acciones de NVIDIA antes del lanzamiento»}: también poco probable. V3 se publicó el 26 de diciembre: ¿quién elegiría el primer día después de Navidad para coordinar una venta en corto? \textit{«Creo que solo estaban apurados por terminar: a quién le importa la Navidad, con tal de publicarlo antes del Año Nuevo chino.»} (Nathan)

\textbf{«Solo costó 6 millones de dólares entrenarlo»}: esto es una lectura errónea grave del artículo. \$5.576M es el costo de horas de GPU de \textbf{una sola ejecución} del preentrenamiento de V3. No incluye los experimentos previos (posiblemente decenas de ablaciones a pequeña escala), el post-entrenamiento, el entrenamiento de R1 ni el servicio de inferencia. La inversión total real podría ser de 10 a 50 veces esa cifra.

\subsection{El Gemini Flash Thinking subestimado}

En las pruebas personales de Lex, el Gemini Flash Thinking de Google podría ser más barato que R1 y no menos capaz, pero casi nadie habla de él. Las razones posibles son:
\begin{itemize}
    \item El marketing de Google es mucho menos llamativo que la estrategia de código abierto de DeepSeek
    \item Flash Thinking probablemente usa un método distinto (agregar razonamiento sobre una arquitectura existente, en lugar de un entrenamiento de razonamiento dedicado)
    \item La narrativa de DeepSeek como «caballo negro chino» se difunde con más facilidad
\end{itemize}

Esta es una lección sobre «tecnología contra narrativa»: la mejor tecnología no necesariamente gana la mayor atención.

\subsection{PyTorch PowerPlant.NoBlowUp}

Al entrenar clústeres grandes, las fluctuaciones de potencia de las GPU pueden destruir el equipo eléctrico. Fase de cómputo con potencia alta (todas las GPU calculando a plena capacidad) $\to$ fase de comunicación con potencia baja (las GPU esperando el intercambio de datos). Esta carga pulsante es muy peligrosa para los transformadores y el equipo de distribución eléctrica.

Un ingeniero de Meta escribió un operador de PyTorch para resolver este problema: durante los periodos de inactividad de las GPU, las hace calcular números sin sentido (multiplicaciones falsas), únicamente para mantener una potencia estable. Este código se hizo de código abierto por accidente: un operador llamado \texttt{PowerPlant.no\_blowup = 1}.

\textit{«Es muy fácil hacer que algo explote.»} (Dylan) Esta historia ilustra perfectamente los retos de ingeniería «aburridos pero cruciales» del entrenamiento de IA: el problema no es un avance algorítmico, sino cómo evitar que tu transformador eléctrico explote.

\subsection{El dilema de seguridad de Anthropic}

Se reporta que Anthropic tiene un modelo de razonamiento mejor que O3, pero no lo publica por consideraciones de seguridad. La cadena de pensamiento (Chain-of-Thought) de R1 en efecto puede resultar inquietante: cambia entre chino e inglés, aparecen fragmentos que parecen texto ininteligible y, de pronto, da la respuesta correcta.

El ritmo agresivo de publicación de DeepSeek redujo los estándares de seguridad de todos, de forma similar a la presión que el programa espacial soviético de la Guerra Fría ejerció sobre la NASA estadounidense. \textit{``DeepSeek ships. That's one of their big advantages.''} (Nathan) La publicación rápida es en sí misma un arma competitiva.

\subsection{La profecía de Sam Altman}

Se cita una frase de Sam Altman que se repite con frecuencia: \textit{«La persuasión sobrehumana llegará antes que la inteligencia sobrehumana.»} Esto tiene implicaciones profundas para la seguridad de la IA: antes de que la IA pueda «pensar de verdad», es posible que ya pueda «persuadir de verdad». Los chatbots de Character AI ya influyen en el estado de ánimo de usuarios jóvenes: ¿qué pasaría si fuera una manipulación cultural o política deliberada?

La declaración pública de Zuckerberg en la llamada de resultados trimestrales también resulta reveladora: \textit{«Para nuestra ventaja nacional, es importante que el estándar de código abierto sea estadounidense.»} El código abierto no es solo una estrategia técnica, sino también una herramienta geopolítica.

\subsection{Resumen de la sección}

La industria de la IA está llena de detalles fascinantes: desde el reto de ingeniería de «no hacer explotar el transformador», hasta la lectura errónea de los «6 millones de dólares», el Gemini subestimado y el código abierto como arma geopolítica. Estas anécdotas revelan una realidad de la industria más rica, más caótica y más humana que la que muestran los artículos técnicos.
\section{Síntesis y ampliación}

\subsection{Teorías de conspiración y verdad}

En torno a DeepSeek circularon teorías de conspiración: ¿el gobierno chino está subsidiando a DeepSeek? Poco probable: la financia un fondo de cobertura, no una empresa vinculada al gobierno. ¿Vendieron en corto acciones de NVIDIA antes del lanzamiento? También poco probable: V3 se lanzó el 26 de diciembre (el primer día después de Navidad). \textit{``Creo que simplemente se apuraron a lanzarlo antes del Año Nuevo chino; a quién le importa la Navidad.''} (Nathan)

Las acciones de NVIDIA se desplomaron: contagio social más una narrativa errónea (``el modelo costó miles de millones''; en realidad ningún modelo público ha costado más de \$1 mil millones de dólares en entrenamiento). La paradoja de Jevons se confirmó pronto: después del lanzamiento de DeepSeek, el precio de las H100 en AWS no bajó sino que subió, y las H200 casi se agotaron.

\subsection{Puntos clave}

\begin{table}[H]
\centering
\begin{tabular}{lp{9cm}}
\toprule
\textbf{Tema} & \textbf{Idea central} \\
\midrule
Tecnología de DeepSeek & MoE + MLA + optimización de bajo nivel en CUDA = ingeniería extrema sobre hardware restringido \\
Controles de exportación & Limitan poco el entrenamiento, pero son más decisivos para limitar el despliegue de inferencia; podrían acelerar la innovación autónoma de China \\
Modelos de razonamiento & R1-Zero es el momento Alpha Zero; los dominios verificables son la clave para escalar el RL \\
TSMC & El mayor punto único de falla de la tecnología mundial; la decadencia de Intel agrava esa concentración \\
Infraestructura & Inversión de \$100 mil millones de dólares a la escala de Stargate; la energía y el contrabando de GPU son las nuevas restricciones \\
Código abierto & La licencia MIT de R1 es un ``gran reinicio''; Tulu de AI2 demuestra que el posentrenamiento abierto puede alcanzar la frontera \\
Panorama de la industria & Solo NVIDIA gana dinero; el problema de los seis nueves de los agentes sigue sin resolverse; la ingeniería de software es donde el cambio es más rápido \\
\bottomrule
\end{tabular}
\end{table}

\subsection{Perspectivas a futuro}

\textbf{Nathan}: la humanidad seguirá existiendo dentro de 1,000 años; no le preocupa que la IA tome el control (razones físicas limitan la velocidad de los robots). Pero le preocupa el \textbf{tecnofascismo}: las interfaces cerebro-máquina crearían una élite fusionada de humanos e IA. Su motivación central: \textit{``No confío en quienes dicen 'confía en mí, hermano, vamos a lograr que la IA sea buena'.''} La IA es la tecnología más poderosa de nuestra vida; se necesita que más personas participen en darle forma.

\textbf{Dylan}: en general, optimista. La IA aumentará la abundancia y reducirá el sufrimiento mediante el mecanismo de la búsqueda de ganancias. Pero le preocupa el pico de sufrimiento humano durante la transición. Lo que más le entusiasma del futuro de la infraestructura: la óptica empaquetada conjuntamente, el entrenamiento en múltiples centros de datos y las nuevas redes de conmutación. \textit{``El ritmo del progreso humano ha alcanzado un nivel sin precedentes.''} Cada capa de la pila de cómputo está innovando: cobre, aire acondicionado, transformadores, litografía.

\textit{``Existe una bondad humana fundamental, y lo que tenemos que hacer es amplificarla.''} (Nathan)

\textit{``Revisar el teléfono mantiene el estado actual del mundo; eso ya es un resultado positivo.''} (Dylan)

\subsection{Lecturas adicionales}

\begin{itemize}
    \item Informe técnico de DeepSeek V3 (diciembre de 2024)
    \item Informe técnico de DeepSeek R1 (enero de 2025)
    \item SemiAnalysis: \href{https://semianalysis.com}{semianalysis.com} — análisis de semiconductores de Dylan Patel
    \item Nathan Lambert: \href{https://www.interconnects.ai}{blog Interconnects} — investigación en IA y código abierto
    \item Proyecto Tulu de AI2: receta de posentrenamiento completamente abierta
    \item Rich Sutton, \textit{The Bitter Lesson}, 2019
    \item Discusión de Andrej Karpathy en Twitter sobre R1-Zero
\end{itemize}

\end{document}
